\section{The \textit{xiéshēng} dogma}
\paragraph{} We have come far enough to articulate the `\textit{xiéshēng} dogma' that will guide our exploration of Old Chinese from here on out. This dogma may now read ``the Old Chinese readings of characters with the same phonetic component share the same rime and a homorganic initial, without mixing stops and resonants''.\footnote{Duan Yucai 段玉裁 (1735-1815) proposed that characters sharing the same phonetic component would have rhymed in the language of the \textit{Shījīng}. Li Fang-Kuei 
\citeyearpar[228]{Li1974} further stipulated that they should have homorganic initials. 
See also \citet[348]{Baxter1992} and \citet[11]{schuessler_minimal_2009}.}
If we come up with a series of sound changes that, if undone reverse chronologically, took us from middle Chinese to a form of Old Chinese that adhered to this dogma, we would have won at the game of Old Chinese reconstruction. 

Once again allowing ourselves to sneak a peak into our own future results, the \emph{xiéshēng} series built on 皮  \emph{bie} %(18-16) 
illustrates
both the rime and the initial facets of the \emph{xiéshēng} hypothesis, in the reconstruction of Baxter and Sagart.

\begin{itemize}
\item {\Large{皮}} \textit{bie} < *{\sil{m-p(r)aj}} `skin`
\item {\Large{被}} 
\textit{bieX} < *{\sil{m-pʰ(r)ajʔ}} `cover`
%\item {\Large{被}} *{\sil{mә-pʰaj}} `cover`
%\item {\Large{被}} *{\sil{pʰajs}} `cover`
\item {\Large{陂}}  \textit{pie}  < *{\sil{p(r)a[j]}} `pond`
\item {\Large{披}} \textit{phie} < *{\sil{pʰ(r)aj}} `drape`
\item {\Large{跛}} \textit{paX} < *{\sil{pˤajʔ}} `lame`
\item {\Large{破}} \textit{phaH} < *{\sil{pʰˤaj-s}} `defeat`
\end{itemize}

\paragraph{Limits of the \textit{xiéshēng} dogma.}
Some series violate the \textit{xiéshēng} dogma, even according to the understanding of those   researchers, such as Baxter and Sagart, who are most thoroughgoing in bringing their reconstructions into conformity with this principle. Witness for example the series built on 
萬 % 21-26/0267a
%    {\large{man}},
     \textit{miənH} `ten thousand'.\footnote{Baxter (personal communication) particularly likes the example of 圓  \textit{ḫiwen} <   *ɢʷ<r>en  `circle', originally written as a circle and with the main vowel *-e-, serving as phonetic in 袁 \textit{ḫiwən} < *[ɢ]ʷa[n] `long' with the main vowel *-a-, but I find this case a bit too complicated to present here.} 


\begin{itemize}
  \item {\Large{萬}} % 21-26/0267a
%    {\large{man}},
     \textit{miənH}
     < *t.ma[n]-s
     `ten thousand' % 0267a

%  \item {\Large{勱}} % 21-26/0267c
    %{\large{{\textsuperscript{virt}}man}},
%    \textit{maeyH} `to strive' % 0267c

  \item {\Large{邁}} % 21-26/0267d
    %{\large{{\textsuperscript{amb}}man}},
    \textit{maeyH} < *mˁrat-s `proceed' 
    
    % 0267d

  \item {\Large{蠆}} % 21-26/0326a
    %{\large{{\textsuperscript{serp}}man}},
    \textit{ṭhaeyH}  < *mә-r̥ˁa[t]-s `scorpion' % 0326a

  \item {\Large{厲}} % 21-26/0340a
    %{\large{{\textsuperscript{cliv}}man}},
    \textit{lieyH} < *[r]at-s `cruel' % 0340a

\begin{itemize}
  \item {\Large{礪}} % 21-26/0340b
    %{\large{{\textsuperscript{lap}}man}},
    \textit{lieyH} < *[r]at-s `grindstone' % 0340b

%  \item {\Large{勵}} % 21-26/0340c
    %{\large{{\textsuperscript{virt}}man}},
%    \textit{lieyH} `to encourage, urge' % 0340c

  \item {\Large{癘}} % 21-26/0340d
    %{\large{{\textsuperscript{infirm}}man}},
    \textit{lieyH} < *[r]at-s `pestilence' % 0340d

  %\item {\Large{\symbol{"275FD}
%}} % 21-26/0340e
    %{\large{{\textsuperscript{insect}}lieyH}},
    %\textit{lieyH} `(rare insect form)' % 0340e

  \item {\Large{蠣}} % 21-26/0340f
    %{\large{{\textsuperscript{serp}}man}},
    \textit{lieyH} < *mә-rat-s `stinging insect' % 0340f

  \item {\Large{糲}} % 21-26/0340g
    %{\large{{\textsuperscript{oryz}}man}},
    \textit{layH} < *[r]ˁat-s, \textit{lat} < *[r]ˁat, \textit{lieyH} 
    < *[r]at-s
    `coarse rice' % 0340g
\end{itemize}
\end{itemize}

Although the Baxter and Sagart reconstructions bring such disparate Middle Chinese readings as \textit{lieyH}, \textit{ṭhaeyH}, and \textit{maeyH} under an overall \textbf{rat}, the base character 萬 \textit{miənH} < *t.ma[n]-s `ten thousand' is precisely missing the element *r- that is otherwise the initial of the series, and shows a final *-n in place of the usual  final *-t. It is not hard to come up with a plausible just so story to account for such a case--the character 萬 originally portrayed a scorpion and thus represented the morpheme  \textit{ṭhaeyH}  < *mә-r̥ˁa[t]-s `scorpion'. As such, it is an excellent base character for the series overall.  In the earliest days of the writing system it was inevitable that characters with more concrete meaning would be pressed into service to represent words with relatively abstract meaning, and needing to represent the word  \textit{miənH}
     < *t.ma[n]-s
     `ten thousand' the ancients turned \textit{faute de mieux} to 萬. 

Let us consider another poorly behaved series, viz. that built on 
柬 \textit{keanX} `select'. % 0185a
This series is uniformly composed of type A readings, but includes the two Old Chinese rimes *-en and *-an. 
The subseries built on 闌 \textit{lan} < *[r]ˁan `entry barrier' specifies  rime  -\textit{an}. 


\begin{itemize}
\item {\Large{柬}} % 23-07/0185a
%{\large{{\textsuperscript{arb}}kean}},
\textit{keanX} < *kˁr[a]nʔ `select' % 0185a

\item {\Large{諫}} % 23-07/0185b
%{\large{{\textsuperscript{dic}}kaen}},
\textit{kaenH} < *kˁranʔ(-s) `admonish' % 0185b

\item {\Large{揀}} % 23-07/0185e
%{\large{{\textsuperscript{man}}kean}},
\textit{keanX} < *kˁr[a]nʔ `select' % 0185e


%\item {\Large{湅}} % 23-07/0185h
%{\large{{\textsuperscript{aq}}len}},
%\textit{lenH} `degum' (silk) % 0185h

\item {\Large{練}} % 23-07/0185i
%{\large{{\textsuperscript{ser}}len}},
\textit{lenH} < *[r]ˁen-s `sized white silk' 

% 0185i

\item {\Large{鍊}} % 23-07/0185j
%{\large{{\textsuperscript{met}}len}},
\textit{lenH} < *[r]ˁen-s `smelt' % 0185j


\item {\Large{闌}} % 23-07/0185f
%{\large{{\textsuperscript{port}}lan}},
\textit{lan} < *[r]ˁan `entry barrier' % 0185f

\begin{itemize}

%\item {\Large{瀾}} % 23-07/0185k
%{\large{{\textsuperscript{aq}}lan}},
%\textit{lan, lanH} `billowing waves'  % 0185k

\item {\Large{瀾}} % 23-07/0185k
%{\large{{\textsuperscript{aq}}lan}},
\textit{lanX} < *[r]ˁanʔ `water of washed rice'  % 0185k

\item {\Large{爛}} % 23-07/0185l
%{\large{{\textsuperscript{ign}}lan}},
\textit{lanH} < *[r]ˁan-s `spoil' % 0185l

%\item {\Large{欄}} % 23-07/0185q
%{\large{{\textsuperscript{arb}}len}},
%\textit{lenH} `Chinaberry tree' % 0185q

\item {\Large{蘭}} % 23-07/0185n
%{\large{{\textsuperscript{herb}}lan}},
\textit{lan} < *k.rˁan `orchid' % 0185n

%\item {\Large{讕}} % 23-07/0185o
%{\large{{\textsuperscript{dic}}lan}},
%\textit{lan, lanX} `slander' % 0185o

\end{itemize}
\end{itemize}

Again we can propose a  plausible just so story. In the earliest days of the writing system, it was inevitable that, simply because insufficient characters were available for every syllable type, scribes could not be too punctilious in their choices. Thus, 柬 was used, \textit{faute de mieux}, for \textbf{krVn}, with an under-specified vowel,  and from there also for both \textbf{ren} and \textbf{ran}. In time, 闌 was specialized for \textbf{ran}.  What I find hard to fathom is why the whole series would be limited to type A rimes. 

I am tempted to draw on Robert Brandom's distinction between a `reason' and a `cause'. The reason why the ancients used 
厲 \textit{lieyH}  `cruel' as phonetic in 礪 \textit{lieyH} `grindstone' is that the morphemes these two characters right have homorganic initials and the same Old Chinese rime.
In contrast, the cause for the ancients reaching for the character 萬 (originally for  \textit{ṭhaeyH}  < *mә-r̥ˁa[t]-s `scorpion') to write the abstract term \textit{miənH} < *t.ma[n]-s `ten thousand' was that they could think of no better option. 
In like fashion, the reason the ancients used 闌 \textit{lan} < *[r]ˁan `entry barrier' as the phonetic in 瀾 \textit{lanX} < *[r]ˁanʔ `water of washed rice' is because of the homorganic initials and shared rime of these two morphemes, but the reason they used 柬 as the phonetic in both 鍊 \textit{lenH} `smelt' and 闌 \textit{lan} `entry barrier' is because, before the coinage of the latter, they had no phonetic available for \textbf{ran}. 

The example of 萬 in particular shows both the power of the \textit{xiéshēng} dogma to make sense of the seeming irregularities in phonetic series and the potential obstacles that face too strict an adherence to this dogma. There is a  tension between the idealized principles of reconstruction and the messy realities of historical usage. 
The methodological question is how far to follow the \textit{xiéshēng}  dogma. I offer the banal contention that one should follow it so long as it is profitable to do so, but no further. I allow myself an analogy from mathematics. Allowing ourselves the operation of division moves us from the integers to the rational numbers, which many will view as worthwhile. However, if we should become so ambitious as to divide by zero, we move instead from the integers to the much simpler universe consisting only of zero itself.\footnote{I remind the reader, that if x/0 = y, then y * 0 = x = 0.} 

The role of the Neogrammarian dogma in historical linguistics is also apt as a comparison. If someone claism that `bridegroom' (cognate with German \textit{Bräutigam}) and `vat' (cognate with \textit{Faß}) disproved the exceptionlessness of sound change, because the first ought to have no -r- and the second ought to start with an f-, this  contention is already undermined by the fact that the mere identification of this contamination (bridegroom) and this borrowing (vat), is only arrived at by prima facie adherence to the principle of exceptionlessness. 
We will only know that we have followed the \textit{xiéshēng} dogma too far when we have in fact followed it too far and seen what consequences we suffer. If we suspect we have gone too far, we can take a step or two backwards if we can determine a plausible `cause' for the  \textit{xiéshēng} dogma's violation in a given instance.

\paragraph{First steps toward the reconstruction of Old Chinese.}
\label{par:first-steps}
To the extent that Old Chinese is more or less like other natural languages, we would expect it to have an inventory of vowels and the vowels of this inventory would more or less each appear as open syllables and preceding the available final consonants that were available in Old Chinese. This is just to say that a rime in Old Chinese was characterized by a vowel and a final (possibly null). We are already operating with the assumption that the two rimes that constitute a \textit{xiéshēng} rime pair characterize a single Old Chinese rime. 

The reconstruction of Old Chinese via \textit{xiéshēng} series can then proceed in two stages. First, we take an inventory of those  \textit{xiéshēng} rime pairs that occur in each syllable type, where a syllable type is a template of a particular initial and a particular final consonant (possibly null). Second, we attempt to identify the \textit{xiéshēng} rime pair pertaining to each distinct syllable type. From the perspective of possible Old Chinese reconstruction, for example, in the first step we identify the the \textit{xiéshēng} rime pairs that characterize Old Chinese *-ek and Old Chinese *-et and in the second step we find a mechanism for identifying that this pair of pairs should be identified as continuing Old Chinese *-e. How to undertake this second step is not immediately clear to me, but I hope it will become so in the course of the work. With which syllable type should be commence our exploration of  \textit{xiéshēng} rime pairs?

To systematically identify  \textit{xiéshēng} rime pairs it behooves us to start with the simplest examples and then look at increasingly complex cases. I have identified syllables with initial 來 \textit{l}- as a good starting point, since this initial is incompatible with quite a few rimes, namely those containing the vowels -\textit{ae}- and -\textit{ea}- (so-called 二等 `division-ii rimes') and those rimes that are only compatible with velars and laryngeals, i.e. % \emph{táng / yáng} 
-\emph{waṅ} (C-I-唐-\textsc{strata}-smooth), 
‑ \emph{iwaṅ} (C-III-陽-\textsc{sol}-smooth),  
%\emph{dēng} 
 \emph{-wəṅ} (C-I-登-\textsc{scando}-smooth),   and 
%\emph{xiān} 
-\emph{wen} (C-IV-先-\textsc{prior}-smooth).\footnote{The two rimes 
%\emph{gēng-èr} 
-\emph{waeṅ} (C-II-庚-\textsc{aetas}-smooth) and  %\emph{gēng} 
 -\emph{weaṅ} (C-II-耕-\textsc{colo}-smooth) belong to both types.} 
